% ============================================================
%  Introducing Python, chapter 2: Types and Variables
% ============================================================
\chapter{Types and Variables}

This chapter explains how a computer stores data, and then how Python handles it in
a way that differs from many other languages. The big idea: in Python a variable is
just a \emph{name} that points at an \emph{object}, and the object knows its own type.
Understanding this now saves you from a whole class of confusing bugs later.

% ------------------------------------------------------------
\section{A Computer}

A program is code plus data. The hardware that holds and processes both has a few parts.

%MD
%```mermaid
%flowchart LR
%    D["Disk: big, slow, keeps data when power is off"] <--> R["RAM: smaller, fast, forgets at power-off"]
%    R <--> C["Cache: tiny, very fast"]
%    C <--> P["CPU: runs instructions"]
%    I["Input: keyboard, mic, camera"] --> P
%    P --> O["Output: screen, speakers"]
%    N["Network"] <--> P
%```
%ENDMD
\begin{center}
\begin{tikzpicture}[font=\sffamily\small,
  n/.style={draw=cblue,fill=cblue!8,rounded corners=2pt,inner sep=5pt,align=center,minimum height=10mm},
  >=Stealth]
  \node[n,draw=cgray,fill=cgray!8] (d) at (0,0) {disk\\[-2pt]{\scriptsize big, slow}};
  \node[n,draw=cteal,fill=cteal!8] (r) at (3.4,0) {RAM\\[-2pt]{\scriptsize fast, limited}};
  \node[n,draw=cgold,fill=cgold!10] (c) at (6.6,0) {cache\\[-2pt]{\scriptsize tiny, faster}};
  \node[n,draw=cred,fill=cred!8] (p) at (9.6,0) {CPU\\[-2pt]{\scriptsize does the work}};
  \draw[<->,thick,cgray] (d) -- (r);
  \draw[<->,thick,cgray] (r) -- (c);
  \draw[<->,thick,cgray] (c) -- (p);
  \node[n,draw=cpurple,fill=cpurple!8] (net) at (5.6,-1.9) {network};
  \node[n,draw=cgreen,fill=cgreen!8] (i) at (8.3,-1.9) {input devices};
  \node[n,draw=cgreen,fill=cgreen!8] (o) at (11.1,-1.9) {output devices};
  \draw[<->,thick,cgray] (net.north) -- (p.south west);
  \draw[->,thick,cgray] (i.north) -- (p.south);
  \draw[->,thick,cgray] (p.south east) -- (o.north);
  \node[font=\scriptsize,text=cgray] at (0,-0.95) {keeps data when off};
  \draw[decorate,decoration={brace,amplitude=4pt,mirror},cgray] (2.4,0.9) -- (10.8,0.9);
  \node[font=\scriptsize,text=cgray] at (6.6,1.35) {volatile: RAM, cache and CPU forget when power is off};
\end{tikzpicture}
\end{center}

Programs and data live on the disk. The CPU can only work on what has been copied into
RAM, so data moves \texttt{disk} $\to$ \texttt{RAM} $\to$ \texttt{cache} $\to$ \texttt{CPU} and back.
A machine learning model works the same way: its weights sit on disk until the program
loads them into memory.

\begin{center}
\begin{tabularx}{\linewidth}{@{}l X l@{}}
\toprule
\textbf{Part} & \textbf{Job} & \textbf{Survives power-off?} \\
\midrule
CPU & runs instructions (many cores in modern chips) & no \\
Cache & small, faster copy of what the CPU needs next & no \\
RAM & working storage for programs and data & no \\
Disk & long-term storage, much larger, much slower & \cmark \\
Input / output & keyboard, mouse, camera, display, printer & n/a \\
Network & talks to other computers & n/a \\
\bottomrule
\end{tabularx}
\end{center}

% ------------------------------------------------------------
\section{Bits and Bytes}

Everything in a computer is made of \textbf{bits}. A bit holds one of two states,
which you can read as off/on, 0/1 or false/true. Eight bits make a \textbf{byte}, which
has $2^8 = 256$ possible patterns. A single bit or byte means nothing by itself. Meaning
comes from how a program chooses to read the pattern.

\subsection{A byte as a number}

Read a byte in base 2 (binary). Each bit stands for a power of two, like the
ones, tens and hundreds places in decimal.

%MD
%```text
%  weight:   128   64   32   16    8    4    2    1
%           +----+----+----+----+----+----+----+----+
%  bits:    | 0  | 0  | 1  | 0  | 0  | 1  | 0  | 1  |
%           +----+----+----+----+----+----+----+----+
%  value:             32          +    4   +    1   = 37
%```
%ENDMD
\begin{center}
\begin{tikzpicture}[font=\ttfamily\small]
  \foreach \w/\b/\f [count=\i from 0] in {128/0/white,64/0/white,32/1/cgold!45,16/0/white,8/0/white,4/1/cgold!45,2/0/white,1/1/cgold!45} {
    \node[font=\sffamily\scriptsize,text=cgray] at (0.95*\i,0.75) {\w};
    \node[draw=cgray,fill=\f,minimum width=9mm,minimum height=9mm] at (0.95*\i,0) {\b};
  }
  \node[font=\sffamily\small,anchor=west] at (7.9,0) {$32+4+1=37$};
  \node[font=\sffamily\scriptsize,text=cgray,anchor=east] at (-0.6,0.75) {weight};
  \node[font=\sffamily\scriptsize,text=cgray,anchor=east] at (-0.6,0) {bits};
\end{tikzpicture}
\end{center}

All eight bits set gives $128+64+32+16+8+4+2+1 = 255$. All clear gives 0.
That is the range of an \emph{unsigned} byte: 0 to 255.

\begin{pycode}
>>> 2 ** 8
256
>>> 0b11111111
255
>>> bin(37)
'0b100101'
>>> f"{37:08b}"
'00100101'
>>> int("00100101", 2)
37
\end{pycode}

A \emph{signed} byte spends one bit on the sign, so it covers $-128$ to $127$ instead.
The same eight bits can therefore mean two different numbers:

\begin{pycode}
>>> (255).to_bytes(1, "big")
b'\xff'
>>> int.from_bytes(b"\xff", "big")
255
>>> int.from_bytes(b"\xff", "big", signed=True)
-1
\end{pycode}

\subsection{A byte as a character}

A byte can also stand for a character. ASCII was the old agreement: 128 codes for English letters,
digits, punctuation and a few control codes. One byte cannot hold every writing
system, so Unicode defines millions of symbols (Chapter 17). Python text uses Unicode.

\begin{pycode}
>>> ord("A")
65
>>> chr(97)
'a'
\end{pycode}

\section{Multibyte Types}

Bytes sit side by side in memory and can be grouped. Two bytes give
$2^{16} = 65{,}536$ patterns; four give $2^{32} \approx 4.3$ billion. Memory locations
are anonymous, so the program has to remember how many bytes belong together
and \emph{what type} they represent.

\begin{pycode}
>>> 2 ** 16
65536
>>> 2 ** 32
4294967296
>>> (1_000_000).to_bytes(4, "big").hex()
'000f4240'
\end{pycode}

%MD
%```text
%  1_000_000 in four bytes (hex):
%  +------+------+------+------+
%  |  00  |  0f  |  42  |  40  |
%  +------+------+------+------+
%   most significant     least significant
%```
%ENDMD
\begin{center}
\begin{tikzpicture}[font=\ttfamily\small]
  \foreach \h/\i in {00/0,0f/1,42/2,40/3}
    \node[draw=cblue,fill=cblue!8,minimum width=14mm,minimum height=9mm] at (1.5*\i,0) {\h};
  \node[font=\sffamily\scriptsize,text=cgray,anchor=east] at (-1,0) {1\,000\,000 =};
  \node[font=\sffamily\scriptsize,text=cgray] at (0,-0.75) {most significant};
  \node[font=\sffamily\scriptsize,text=cgray] at (4.5,-0.75) {least significant};
\end{tikzpicture}
\end{center}

\begin{keyidea}[: the same bytes, different meanings]
Bytes carry no meaning until a type is applied. The two bytes \code{41 42} are the text
\code{AB}, the integer 16706, or the list of small numbers 65 and 66, depending on the type
you read them as.
\end{keyidea}

%MD
%```mermaid
%flowchart LR
%    B["bytes: 41 42"] -->|read as text| T["'AB'"]
%    B -->|read as one integer| I["16706"]
%    B -->|read as separate numbers| L["[65, 66]"]
%```
%ENDMD
\begin{center}
\begin{tikzpicture}[font=\sffamily\small,>=Stealth,
  o/.style={draw=cgreen,fill=cgreen!10,rounded corners=2pt,inner sep=4pt,font=\ttfamily\small}]
  \node[draw=cblue,fill=cblue!8,font=\ttfamily\small,inner sep=5pt] (b) at (0,0) {41 42};
  \node[o] (t) at (5,1) {'AB'};
  \node[o] (i) at (5,0) {16706};
  \node[o] (l) at (5,-1) {[65, 66]};
  \draw[->,thick,cgray] (b) -- node[above,sloped,font=\scriptsize]{as text} (t);
  \draw[->,thick,cgray] (b) -- node[above,font=\scriptsize]{as one integer} (i);
  \draw[->,thick,cgray] (b) -- node[below,sloped,font=\scriptsize]{as separate numbers} (l);
\end{tikzpicture}
\end{center}

\begin{pycode}
>>> data = b"AB"
>>> data.decode("ascii")
'AB'
>>> int.from_bytes(data, "big")
16706
>>> list(data)
[65, 66]
\end{pycode}

% ------------------------------------------------------------
\section{Variables}

A \textbf{variable} is a name in your code that refers to a value. Python handles this
differently from languages like C or Java, so the picture matters.

\subsection{Assign a value}

\begin{pycode}
>>> temperature = 0.7
>>> temperature
0.7
\end{pycode}

The shape is \code{name = value}. The \code{=} sign means \emph{assign}, not ``equals''.
The right side is evaluated first; the result is then given the name on the left.
Typing a bare name in the interpreter prints its value. A file would not.

\subsection{A name is a sticky note}

In many languages a variable is a \emph{box} with a fixed type, and assignment copies the
value into the box. In Python, the value lives in an \textbf{object} somewhere in memory.
The variable is only a \textbf{name} (also called a reference or label) attached to that object.
Python is \textbf{dynamically typed}: the type belongs to the object, not to the name.

%MD
%```text
%  Other languages:             Python:
%
%   x  +-----+                   x ------> +--------------+
%      |  5  |   (box holds        name     | int  |  5    |   (object holds
%      +-----+    the value)                +--------------+    value + type)
%```
%ENDMD
\begin{center}
\begin{tikzpicture}[font=\sffamily\small,>=Stealth,
  nm/.style={draw=cgold,fill=cgold!25,rounded corners=1pt,inner sep=3pt,font=\ttfamily\small},
  ob/.style={draw=cblue,fill=cblue!8,minimum height=8mm,inner sep=4pt}]
  \node[font=\sffamily\bfseries\small] at (0.8,1.4) {Box model (C, Java)};
  \node[nm] at (0,0) {x};
  \node[draw=cgray,fill=cgray!8,minimum width=12mm,minimum height=8mm,font=\ttfamily] at (1.7,0) {5};
  \node[font=\scriptsize,text=cgray] at (0.9,-0.8) {the name \emph{is} the box};
  \node[font=\sffamily\bfseries\small] at (8.2,1.4) {Python};
  \node[nm] (x) at (6.4,0) {x};
  \node[ob] (o) at (9.6,0) {\ttfamily int \ \ 5};
  \draw[->,thick] (x) -- (o);
  \node[font=\scriptsize,text=cgray] at (8.2,-0.8) {the name points at an object};
\end{tikzpicture}
\end{center}

\begin{pycode}
>>> x = 5
>>> type(x)
<class 'int'>
>>> x = "five"
>>> type(x)
<class 'str'>
\end{pycode}

The name \code{x} was reused for a different type without trouble. Only the object
\code{5} is an \code{int}; only \code{"five"} is a \code{str}.

\subsection{What an object holds}

Every Python object carries some bookkeeping with its value.

\begin{center}
\begin{tabularx}{\linewidth}{@{}l X@{}}
\toprule
\textbf{Part} & \textbf{Meaning} \\
\midrule
\code{id} & a unique number for this object while it lives; \code{id(obj)} \\
type & what kind of object it is; \code{type(obj)} \\
value & the actual data \\
reference count & how many names (and other holders) point at it \\
\bottomrule
\end{tabularx}
\end{center}

When the reference count drops to zero, nothing can reach the object any more and
Python reclaims its memory. This is \textbf{garbage collection}. You never call
\code{malloc} or \code{free} as in C.

\begin{pycode}
>>> import sys
>>> a = []
>>> sys.getrefcount(a)
2
>>> b = a
>>> sys.getrefcount(a)
3
\end{pycode}

(\code{getrefcount} reports one more than you expect, because passing \code{a} to the function
adds a temporary reference.)

\subsection{Change a value: rebind the name}

Assigning again does \emph{not} edit the old object. Python builds (or finds) another
object and moves the name to it.

%MD
%```mermaid
%flowchart LR
%    subgraph before ["x = 5"]
%    x1((x)) --> o5["int 5"]
%    end
%    subgraph after ["x = 6"]
%    x2((x)) --> o6["int 6"]
%    o5b["int 5 (refs 0, freed)"]
%    end
%```
%ENDMD
\begin{center}
\begin{tikzpicture}[font=\sffamily\small,>=Stealth,
  nm/.style={draw=cgold,fill=cgold!25,rounded corners=1pt,inner sep=3pt,font=\ttfamily\small},
  ob/.style={draw=cblue,fill=cblue!8,minimum height=8mm,inner sep=4pt,font=\ttfamily\small}]
  \node[font=\scriptsize\bfseries,text=cgray] at (2.3,1.2) {before: \texttt{x = 5}};
  \node[nm] (x1) at (0,0) {x};
  \node[ob] (a) at (3.5,0) {int 5};
  \draw[->,thick] (x1) -- (a);
  \node[font=\scriptsize\bfseries,text=cgray] at (9.3,1.2) {after: \texttt{x = 6}};
  \node[nm] (x2) at (7,0) {x};
  \node[ob] (b) at (11,0.5) {int 6};
  \node[ob,draw=cgray,dashed,text=cgray,fill=white] (c) at (11,-0.7) {int 5 (refs 0: freed)};
  \draw[->,thick] (x2) -- (b);
\end{tikzpicture}
\end{center}

\begin{pycode}
>>> x = 1000
>>> first = id(x)
>>> x = 2000
>>> id(x) == first
False
\end{pycode}

\begin{trap}
CPython keeps ready-made objects for small integers ($-5$ to 256) and reuses them. So
\code{is} can seem to work on small numbers and fail on larger ones. Compare values with \code{==}.
Use \code{is} only for \code{None}.
\begin{pycode}
>>> p = int("256"); q = int("256")
>>> p is q
True
>>> p = int("1000"); q = int("1000")
>>> p is q
False
>>> p == q
True
\end{pycode}
\end{trap}

\subsection{Delete a variable}

\code{del name} removes the name (and lowers the object's count). You rarely need it, because
Python frees names when their scope ends.

\begin{pycode}
>>> tmp = 3
>>> del tmp
>>> tmp
Traceback (most recent call last):
  ...
NameError: name 'tmp' is not defined
\end{pycode}

\subsection{Two names, one object}

Because names are labels, \code{b = a} does not copy anything. It attaches a second label to the
same object. For immutable objects (numbers, text) this is harmless. For mutable ones,
it is a classic surprise.

\begin{boss}[: aliasing a mutable object]
Two names, one list. A change made through one name is visible through the other, because
there is only one object.
\medskip
%MD
%```mermaid
%flowchart LR
%    a((a)) --> L["list [1, 2, 3]"]
%    b((b)) --> L
%    c((c)) --> M["list [1, 2, 3] (separate copy)"]
%```
%ENDMD
\begin{center}
\begin{tikzpicture}[font=\sffamily\small,>=Stealth,
  nm/.style={draw=cgold,fill=cgold!25,rounded corners=1pt,inner sep=3pt,font=\ttfamily\small},
  ob/.style={draw=cblue,fill=cblue!8,minimum height=8mm,inner sep=4pt,font=\ttfamily\small}]
  \node[nm] (a) at (0,0.6) {a};
  \node[nm] (b) at (0,-0.6) {b};
  \node[ob] (l) at (3.6,0) {list [1, 2, 3]};
  \draw[->,thick] (a) -- (l);
  \draw[->,thick] (b) -- (l);
  \node[nm] (c) at (7,0) {c};
  \node[ob] (m) at (10.6,0) {list [1, 2, 3]};
  \draw[->,thick] (c) -- (m);
  \node[font=\scriptsize,text=cgray] at (1.8,-1.3) {\texttt{b = a}: one object};
  \node[font=\scriptsize,text=cgray] at (8.8,-1.3) {\texttt{c = a.copy()}: a second object};
\end{tikzpicture}
\end{center}
\begin{pycode}
>>> a = [1, 2]
>>> b = a
>>> b.append(3)
>>> a
[1, 2, 3]
>>> a is b
True
>>> c = a.copy()
>>> c.append(4)
>>> a, c
([1, 2, 3], [1, 2, 3, 4])
\end{pycode}
Ask yourself: ``how many objects exist here?'' If you want an independent copy, make one
on purpose (\code{.copy()}); lists are covered in Chapter 8.
\end{boss}

\subsection{Name variables}

Names follow hard rules and soft conventions.

\begin{center}
\begin{tabularx}{\linewidth}{@{}l X X@{}}
\toprule
\textbf{Rule (enforced)} & \textbf{Valid} & \textbf{Invalid} \\
\midrule
letters, digits, underscore only & \code{max\_tokens2} & \code{max-tokens}, \code{name!} \\
may not start with a digit & \code{\_1st}, \code{first\_1} & \code{1st\_place} \\
case-sensitive (these are three names) & \code{Model}, \code{model}, \code{MODEL} & --- \\
not a reserved word & \code{class\_name} & \code{class}, \code{for}, \code{None} \\
\bottomrule
\end{tabularx}
\end{center}

\begin{pycode}
>>> 1st_place = 1
Traceback (most recent call last):
  ...
SyntaxError: invalid decimal literal
>>> my name = 5
Traceback (most recent call last):
  ...
SyntaxError: invalid syntax
>>> import keyword
>>> len(keyword.kwlist)
35
>>> keyword.iskeyword("lambda")
True
\end{pycode}

\begin{center}
\begin{tabularx}{\linewidth}{@{}l l X@{}}
\toprule
\textbf{Style} & \textbf{Example} & \textbf{Use it for} \\
\midrule
snake\_case & \code{num\_tokens} & variables, functions, modules \\
PascalCase & \code{ChatClient} & classes (Chapter 11) \\
UPPER\_SNAKE & \code{MAX\_RETRIES} & ``constants'': Python does not enforce them \\
leading underscore & \code{\_cache} & ``internal, please do not import'' \\
double underscores both sides & \code{\_\_init\_\_} & special (``dunder'') methods \\
camelCase & \code{numTokens} & not Python style (JavaScript, Java) \\
\bottomrule
\end{tabularx}
\end{center}

\begin{gamerule}
Name things for what they hold: \code{prompt\_tokens}, not \code{x1}. Never use a built-in as a name
(\code{list}, \code{str}, \code{id}, \code{type}); it hides the real one.
\end{gamerule}

% ------------------------------------------------------------
\section{Python Types}

Bits are meaningless until a type reads them. Python provides built-in types that match
common hardware types, plus some software-only ones. The column that matters most for
bugs is \textbf{mutable}: can the value be changed in place after it is created?

\begin{center}
\begin{tabularx}{\linewidth}{@{}l l c l l@{}}
\toprule
\textbf{Name} & \textbf{Type} & \textbf{Mutable?} & \textbf{Example} & \textbf{Chapter} \\
\midrule
Boolean & \code{bool} & no & \code{True}, \code{False} & 3 \\
Integer & \code{int} & no & \code{47}, \code{25\_000} & 3 \\
Float & \code{float} & no & \code{3.14}, \code{2.7e5} & 3 \\
Complex & \code{complex} & no & \code{3j}, \code{5 + 9j} & 3 \\
Text string & \code{str} & no & \code{"alas"}, \code{'alack'} & 4 \\
Bytes & \code{bytes} & no & \code{b'ab\textbackslash xff'} & 5 \\
Bytearray & \code{bytearray} & \textbf{yes} & \code{bytearray(b'ab')} & 5 \\
Tuple & \code{tuple} & no & \code{(2, 4, 8)} & 8 \\
List & \code{list} & \textbf{yes} & \code{["a", "b"]} & 8 \\
Set & \code{set} & \textbf{yes} & \code{\{3, 5, 7\}} & 9 \\
Frozen set & \code{frozenset} & no & \code{frozenset([1, 2])} & 9 \\
Dictionary & \code{dict} & \textbf{yes} & \code{\{"a": 1\}} & 9 \\
Nothing & \code{NoneType} & no & \code{None} & here \\
\bottomrule
\end{tabularx}
\end{center}

The last row is not in the table of the book, but you will use it constantly: \code{None}
means ``no value here''.

Two functions let you ask an object about its type:

\begin{pycode}
>>> type(3.5)
<class 'float'>
>>> isinstance(3.5, float)
True
>>> isinstance(True, int)
True
\end{pycode}

(\code{bool} is a kind of \code{int} in Python, which Chapter 3 explains.)

\begin{typecheck}[: mutable or immutable?]
\textbf{Immutable} means no operation can change the object in place; you can only
make a new one. \textbf{Mutable} means it can change, and every name pointing at it sees the change.
\medskip
%MD
%```mermaid
%flowchart LR
%    subgraph immutable ["str: immutable"]
%    s1["'abc'"] -->|"s.upper()"| s2["'ABC' (new object)"]
%    end
%    subgraph mutable ["list: mutable"]
%    l1["[1, 2]"] -->|"l.append(3)"| l2["same object, now [1, 2, 3]"]
%    end
%```
%ENDMD
\begin{center}
\begin{tikzpicture}[font=\ttfamily\small,>=Stealth,
  ob/.style={draw=cblue,fill=cblue!8,minimum height=7mm,inner sep=4pt}]
  \node[font=\sffamily\scriptsize\bfseries,text=cgray] at (2.3,1.2) {str: immutable};
  \node[ob] (s1) at (0,0) {"abc"};
  \node[ob] (s2) at (4.6,0) {"ABC"};
  \draw[->,thick,cgray] (s1) -- node[above,font=\scriptsize]{s.upper()} node[below,font=\sffamily\scriptsize]{new object} (s2);
  \node[font=\sffamily\scriptsize\bfseries,text=cgray] at (9.9,1.2) {list: mutable};
  \node[ob] (l1) at (7.6,0) {[1, 2]};
  \node[ob] (l2) at (12.2,0) {[1, 2, 3]};
  \draw[->,thick,cgray] (l1) -- node[above,font=\scriptsize]{l.append(3)} node[below,font=\sffamily\scriptsize]{same object} (l2);
\end{tikzpicture}
\end{center}
\begin{pycode}
>>> s = "abc"
>>> s[0] = "X"
Traceback (most recent call last):
  ...
TypeError: 'str' object does not support item assignment
>>> items = ["a", "b", "c"]
>>> before = id(items)
>>> items[0] = "X"
>>> items
['X', 'b', 'c']
>>> id(items) == before
True
\end{pycode}
\end{typecheck}

\begin{typecheck}[: \texttt{None} and truthiness]
\code{None} is the only value of type \code{NoneType}. It is falsy. Test for it with
\code{is None}, and never call methods on a value that might be \code{None}.
\begin{pycode}
>>> result = None
>>> result is None
True
>>> bool(result)
False
>>> result.upper()
Traceback (most recent call last):
  ...
AttributeError: 'NoneType' object has no attribute 'upper'
\end{pycode}
\end{typecheck}

\subsection{Type hints do not enforce types}

You may annotate a name with a type. Python records the annotation but does not check it
while running. Tools such as mypy or your editor check it before you run (Chapter 14).

\begin{pycode}
>>> count: int = "five"
>>> count
'five'
\end{pycode}

\section{Specify Values}

A \textbf{literal} is a value written directly in code: \code{5}, \code{"hi"}, \code{[1, 2]}.
A \textbf{variable} is a name that refers to a value. Each type has its own literal syntax,
and punctuation decides the type. Two cases catch beginners:

\begin{center}
\begin{tabularx}{\linewidth}{@{}l l X@{}}
\toprule
\textbf{You write} & \textbf{You get} & \textbf{Why} \\
\midrule
\code{(4)} & \code{int} & parentheses only group; they do not make a tuple \\
\code{(4,)} & \code{tuple} & the comma makes the tuple \\
\code{\{\}} & \code{dict} & empty braces are a dict; use \code{set()} for an empty set \\
\code{\{4\}} & \code{set} & braces with items but no colons \\
\code{\{"a": 4\}} & \code{dict} & braces with \code{key: value} pairs \\
\bottomrule
\end{tabularx}
\end{center}

\begin{pycode}
>>> type((4))
<class 'int'>
>>> type((4,))
<class 'tuple'>
>>> type({})
<class 'dict'>
>>> type({4})
<class 'set'>
\end{pycode}

\section{Objects as Plastic Boxes in Memory}

A good mental image: each object is a clear plastic box on a memory shelf. Some contents
have a fixed size (id, type, count); the value part can be small or large. A
list box can sit beside an int box. Your code lives on another shelf, and each variable is a
sticky note stuck to a box.

%MD
%```text
%  code shelf                    data shelf
%  ----------                    ----------------------------------------
%  prompt  --note-->  [ id | type: str  | "Hi"       | refs: 1 ]
%  n_calls --note-->  [ id | type: int  | 3          | refs: 1 ]
%  history --note-->  [ id | type: list | [..., ...] | refs: 1 ]
%```
%ENDMD
\begin{center}
\begin{tikzpicture}[font=\ttfamily\scriptsize,>=Stealth,
  nm/.style={draw=cgold,fill=cgold!25,rounded corners=1pt,inner sep=3pt},
  c/.style={draw=cblue,fill=cblue!8,minimum height=6mm,inner sep=3pt}]
  \foreach \n/\t/\v/\y [count=\k] in {prompt/str/"Hi"/0.9,n\_calls/int/3/0,history/list/{[..., ...]}/-0.9} {
    \node[nm,anchor=east] (n\k) at (1.6,\y) {\n};
    \node[c] (i\k) at (3.4,\y) {id};
    \node[c,right=0pt of i\k] (t\k) {type: \t};
    \node[c,right=0pt of t\k] (v\k) {\v};
    \node[c,right=0pt of v\k] (r\k) {refs: 1};
    \draw[->,thick] (n\k) -- (i\k);
  }
\end{tikzpicture}
\end{center}

\section{Review}

Variables are names; objects carry an id, a type, a value and a reference count. Assigning
moves a name to another object. Mutable objects can change in place, so shared names can
surprise you. Next: numbers, the simplest objects.

\begin{practice}
\begin{enumerate}[itemsep=1pt]
\item Why can the same name hold an \code{int} and, a moment later, a \code{str}?
\item After \code{a = [1]} and \code{b = a}, then \code{b.append(2)}, what is \code{a}? Why?
\item Which of these is not a valid variable name, and why: \code{top\_p}, \code{2nd}, \code{Top\_P}, \code{for}?
\item What is the difference between a mutable and an immutable type? Name two of each.
\item Why is \code{x is 1000} unreliable, and what should you write instead?
\item What type is \code{(7)}, and what must you add to make it a tuple?
\end{enumerate}
\end{practice}

\begin{achievement}{Name Tamer}
\begin{itemize}[itemsep=1pt]
\item Explain a variable as a name pointing at an object, with a type, id and reference count.
\item Predict what happens to other names when a mutable object is changed.
\item Choose valid, idiomatic names and spot illegal ones.
\item Read the type table and say which types are mutable.
\item Check types with \code{type()} and \code{isinstance()}, and test for \code{None} with \code{is}.
\end{itemize}
\end{achievement}
